\documentclass[a4paper,12pt]{article}

\usepackage[utf8]{inputenc}
\usepackage[T1,T2A]{fontenc}
\usepackage[russian]{babel}
\usepackage{cmap}
\usepackage{dejavu}
\usepackage{geometry}
\usepackage{amsmath,amssymb}
\usepackage[dvipsnames,table]{xcolor}
\usepackage{tikz}
\usetikzlibrary{arrows.meta,positioning,calc,fit,backgrounds,shapes.geometric}
\usepackage[hidelinks]{hyperref}

\geometry{top=25mm,bottom=30mm,left=25mm,right=20mm}
\setlength{\parindent}{0pt}
\setlength{\parskip}{0.55em}
\pagestyle{plain}

\author{Выполнил: Фамилия И.\,О.\\Группа: группа}
\title{Лабораторная работа по TikZ}
\date{\today}

\definecolor{softblue}{RGB}{226,240,255}
\definecolor{softgreen}{RGB}{230,247,236}
\definecolor{softorange}{RGB}{255,243,224}
\definecolor{deepblue}{RGB}{29,78,137}
\definecolor{deepgreen}{RGB}{29,121,76}
\definecolor{deeporange}{RGB}{173,96,0}

\begin{document}
\maketitle

\textbf{Идея работы.} В качестве основы взят пример из файла tikz.pdf из архива \texttt{tikz.zip}. Дополнительно использован пример блок-схемы из block.pdf. Оба исходника не просто перекрашены: блок-схема переработана в более компактную и аккуратную композицию, а координатный пример превращён в самостоятельный графический сюжет с осями, областями, траекториями и подписями.

\section*{1. Переработанная блок-схема}

\begin{center}
\begin{tikzpicture}[
    >=Stealth,
    node distance=11mm and 16mm,
    every node/.style={font=\small},
    start/.style={ellipse, draw=deepblue, fill=softblue, thick, minimum width=31mm, minimum height=9mm, align=center, font=\bfseries},
    process/.style={rectangle, rounded corners=2.5pt, draw=deepgreen, fill=softgreen, thick, minimum width=43mm, minimum height=10mm, align=center},
    decision/.style={diamond, aspect=2.0, draw=deeporange, fill=softorange, thick, inner sep=1pt, text width=28mm, align=center},
    arrow/.style={->, thick, draw=black!80}
]
    \node[start] (begin) {Начало\\входных данных};
    \node[process, below=of begin] (read) {Считывание и проверка структуры};
    \node[decision, below=of read, yshift=-1mm] (ok1) {Всё\\корректно?};
    \node[process, below left=13mm and 24mm of ok1] (fix) {Исправление\\ошибок};
    \node[process, below=of ok1] (normalize) {Нормализация\\параметров};
    \node[process, below=of normalize] (solve) {Вычисление\\результата};
    \node[decision, below=of solve, yshift=-1mm] (ok2) {Ответ\\подходит?};
    \node[process, below left=13mm and 24mm of ok2] (retune) {Уточнение\\условий};
    \node[process, below=of ok2] (report) {Формирование\\отчёта};
    \node[start, below=of report] (finish) {Готово};

    \draw[arrow] (begin) -- (read);
    \draw[arrow] (read) -- (ok1);
    \draw[arrow] (ok1) -- node[right]{да} (normalize);
    \draw[arrow] (ok1.west) -- ++(-3.2,0) |- node[pos=0.25,left]{нет} (fix.north);
    \draw[arrow] (fix.east) -| (read.west);

    \draw[arrow] (normalize) -- (solve);
    \draw[arrow] (solve) -- (ok2);
    \draw[arrow] (ok2) -- node[right]{да} (report);
    \draw[arrow] (ok2.west) -- ++(-3.2,0) |- node[pos=0.25,left]{нет} (retune.north);
    \draw[arrow] (retune.east) -| (solve.west);

    \draw[arrow] (report) -- (finish);

    \begin{scope}[on background layer]
        \node[rounded corners=4mm, draw=deepblue!25, fill=deepblue!4, inner sep=7mm, fit=(begin)(read)(ok1)(fix), label={[font=\bfseries\small]above:Подготовка}] {};
        \node[rounded corners=4mm, draw=deepgreen!25, fill=deepgreen!4, inner sep=7mm, fit=(normalize)(solve)(ok2)(retune)(report)(finish), label={[font=\bfseries\small]above:Обработка и вывод}] {};
    \end{scope}
\end{tikzpicture}
\end{center}

\clearpage
\section*{2. Переработанный координатный рисунок}

\begin{center}
\begin{tikzpicture}[scale=0.92, every node/.style={font=\small}]
    \fill[white] (-0.8,-0.8) rectangle (10.8,8.8);
    \draw[step=1, very thin, gray!18] (0,0) grid (10,8);
    \draw[->, thick] (0,0) -- (10.3,0) node[right] {$x$};
    \draw[->, thick] (0,0) -- (0,8.4) node[above] {$y$};

    \fill[Blue!10, draw=Blue!30] (1.0,6.7) .. controls (2.0,7.8) and (3.6,7.7) .. (4.9,6.7)
        .. controls (6.1,5.8) and (7.1,5.1) .. (8.7,5.2)
        .. controls (9.3,5.2) and (9.6,4.6) .. (9.2,4.0)
        .. controls (8.7,3.0) and (7.5,2.4) .. (6.2,2.5)
        .. controls (4.3,2.7) and (2.0,4.0) .. (1.0,6.7) -- cycle;

    \draw[very thick, black!85] (0.4,7.7) .. controls (2.3,7.1) and (3.7,6.0) .. (5.0,4.8)
        .. controls (6.2,3.7) and (7.7,2.6) .. (9.6,1.7);

    \draw[very thick, deepgreen] (1.0,1.0) .. controls (2.6,1.8) and (4.2,3.0) .. (5.0,4.2)
        .. controls (6.0,5.7) and (7.1,6.6) .. (9.5,7.1);

    \draw[thick, dashed, deeporange] (0.8,1.0) -- (2.1,1.9) -- (3.2,2.7) -- (4.3,3.6) -- (5.2,4.6) -- (6.5,5.7) -- (8.1,6.5);

    \draw[very thick, Violet] (5.2,0.2) -- (5.2,8.2);

    \foreach \p/\lab in {
      (1.2,6.9)/A,
      (2.3,5.2)/B,
      (3.7,6.1)/C,
      (5.2,4.3)/D,
      (6.4,6.6)/E,
      (7.3,4.5)/F,
      (8.2,3.0)/G,
      (9.0,5.6)/H,
      (2.0,2.1)/I,
      (3.0,1.5)/J,
      (6.1,2.2)/K,
      (7.8,1.2)/L
    }{%
        \fill[MidnightBlue] \p circle (2.4pt);
        \node[font=\bfseries\small, text=MidnightBlue, xshift=1.7mm, yshift=1.7mm] at \p {$\lab$};
    }

    \node[fill=white, fill opacity=0.9, text opacity=1, draw=black!25, rounded corners=2pt, inner sep=2pt, anchor=north west] at (0.25,7.95) {$\mathcal{D}$};
    \node[fill=white, fill opacity=0.9, text opacity=1, draw=black!25, rounded corners=2pt, inner sep=2pt, anchor=south east] at (9.85,0.15) {Координатный сюжет};
\end{tikzpicture}
\end{center}

\end{document}
